\documentclass[conference]{IEEEtran}
\IEEEoverridecommandlockouts

\usepackage[english]{babel}
\usepackage{amsmath,amssymb}
\usepackage{graphicx}
\usepackage{textcomp}
\usepackage{booktabs}
\usepackage{url}
\usepackage[hidelinks]{hyperref}
\usepackage{cite}
\usepackage{array}
\usepackage{microtype}

\title{Enhancing Graph Neural Networks with Retrieval-Augmented Generation for Node Classification}

\author{%
\IEEEauthorblockN{Neeraj Yadav\textsuperscript{*},
Dev Singh\textsuperscript{**},
Gourav Varanasi}%
\IEEEauthorblockA{\textit{Department of Computer Science}\\
\textit{Shiv Nadar University}\\
Greater Noida, India\\
\textsuperscript{*}Student ID: 2310110199 \quad
\textsuperscript{**}Student ID: 2310110099}%
}

\begin{document}
\maketitle

\begin{abstract}
We present a social network intelligence system that combines Graph Neural Networks (GNNs) for structure-aware prediction, hybrid retrieval over a Neo4j graph database with native vector indexes, and a multi-agent Knowledge-Augmented Generation (KAG) pipeline driven by a large language model. The work progressed from exploratory RAG-only notebooks to GNN baselines, then to a GNN\,+\,semantic feature hybrid trained on the SNAP Facebook page--page network, and finally to an inference-time graph-RAG evaluation with validation-only hyperparameter search and automated regression tests. This report summarizes motivation, literature context, system architecture, Neo4j-backed retrieval, training on Kaggle, and the test suite used to validate pipeline accuracy.
\end{abstract}

\begin{IEEEkeywords}
Graph neural networks, GraphRAG, Neo4j, vector search, hybrid retrieval, social networks, knowledge-augmented generation, node classification.
\end{IEEEkeywords}

\section{Introduction}

Social graphs encode both \emph{who connects to whom} and rich metadata about users and content. Pure text retrieval ignores topology; pure graph algorithms may miss semantics in unstructured text. Recent surveys argue that combining graph machine learning with large language models (LLMs) is an active frontier: graphs supply relational inductive bias and grounding, while LLMs supply flexible reasoning and natural-language interfaces~\cite{acm_survey_gml_llm}.

This project implements an end-to-end pipeline: (i) GNN-based encoders and classifiers trained on real social datasets; (ii) ingestion of graph and text into Neo4j with vector indexes for similarity search; (iii) a FastAPI service orchestrating analyzer, router, hybrid retrievers, synthesizer, and validator agents; and (iv) automated tests including a stratified node-classification pipeline accuracy gate. The progression from early notebooks to production-style code is documented in the project training log~\cite{project_doc_internal}.

\section{Literature Review: Graph ML and LLMs}

\subsection{ACM Survey on Graph ML and LLMs}

We position our work within the ACM 2025 survey \emph{Graph Machine Learning in the Era of Large Language Models}~\cite{acm_survey_gml_llm}. That survey discusses how LLMs can assist graph tasks (e.g., feature augmentation, explanation, query generation) and how graph structure can reduce hallucination and improve factual grounding for retrieval-augmented generation. Themes relevant to our implementation include:

\begin{itemize}
  \item \textbf{Hybrid modeling:} combining parametric graph models (GNNs) with non-parametric retrieval over documents or subgraphs.
  \item \textbf{Retrieval and memory:} using external stores (Neo4j with vector properties) as scalable memory for LLM prompts.
  \item \textbf{Evaluation rigor:} avoiding data leakage and reporting clear train/validation/test protocols, motivating our three-way split and regression tests.
\end{itemize}

\subsection{Supporting References}

Graph convolutional and sampling approaches (e.g., GraphSAGE) underpin our encoders~\cite{hamilton2017graphsage}. PyTorch Geometric provides scalable GNN primitives~\cite{pyg}. Neo4j supports property graphs and vector indexes for combined structural and semantic queries~\cite{neo4j_vector}.

\section{Problem Setting and Datasets}

We focus on node-level tasks on social graphs---notably the SNAP Facebook Large Page--Page Network (MUSAE-style CSVs: edges, targets, optional features) with four page-type classes. Edges are treated as undirected; features may be sparse binary indicators and/or concatenated sentence embeddings. Table~\ref{tab:data} summarizes the main dataset characteristics.

\begin{table}[t]
\centering
\caption{Facebook page--page network (representative scale).}
\label{tab:data}
\begin{tabular}{@{}l r@{}}
\toprule
Quantity & Approx. value \\
\midrule
Nodes & $\sim$22{,}470 \\
Edges (undirected) & $\sim$171{,}000 \\
Classes & 4 (e.g., politician, government, tvshow, company) \\
\bottomrule
\end{tabular}
\end{table}

\section{Evolution of Methodology}

Following~\cite{project_doc_internal}, development advanced in four stages:

\begin{enumerate}
  \item \textbf{RAG-only (Notebook 1):} nodes converted to text, embedded (e.g., MiniLM), retrieved with FAISS/Gemini---useful for exploration but not a principled node classifier; graph structure unused.
  \item \textbf{GNN-only (Notebook 2):} 2-layer GraphSAGE on sparse features; stratified split; test accuracy $\sim$93.5\%.
  \item \textbf{GNN\,+\,RAG hybrid training (Notebook 3):} semantic embeddings concatenated to features; 3-layer SAGE with BatchNorm/Dropout; reported test accuracy $\sim$94.6\%; checkpoint saved for downstream use.
  \item \textbf{Inference-time graph-RAG + tests:} load checkpoint; multi-hop label propagation and homophily-gated blending tuned on a \emph{validation} slice only; pytest regression with a high accuracy threshold.
\end{enumerate}

\section{System Architecture}

\subsection{Multi-Agent Query Pipeline}

Figure~\ref{fig:arch} depicts the runtime flow: Query Analyzer $\rightarrow$ Router $\rightarrow$ Graph/Vector retrievers (with fusion) $\rightarrow$ GNN inference $\rightarrow$ Synthesizer (KAG) $\rightarrow$ Validator, consistent with the repository README.

\begin{figure}[t]
\centering
\includegraphics[width=0.95\linewidth]{architecture_diagram.png}
\caption{High-level multi-agent pipeline for queries over the social graph.}
\label{fig:arch}
\end{figure}

\subsection{Neo4j as Unified Backend}

The system uses Neo4j for both graph traversal and vector similarity~\cite{arch_md_internal}. Figure~\ref{fig:neo4j} illustrates the graph/vector backend (Neo4j Browser or equivalent view).

\begin{figure}[t]
\centering
\includegraphics[width=0.95\linewidth]{neo4j.png}
\caption{Neo4j property graph and vector-backed retrieval context.}
\label{fig:neo4j}
\end{figure}

Key design points (see \texttt{ARCHITECTURE.md} and \texttt{rag/hybrid\_retrieval.py}):

\begin{itemize}
  \item \textbf{Embeddings on nodes:} text embeddings (e.g., 384-d) and GNN embeddings (e.g., 128-d) stored as node properties; separate vector indexes per dimensionality.
  \item \textbf{Hybrid Cypher:} hot paths combine pattern matching with vector scoring where possible to reduce round-trips.
  \item \textbf{Python RRF:} reciprocal rank fusion when two distinct Cypher patterns must be merged.
\end{itemize}

\section{Implementation Mapping}

Table~\ref{tab:modules} maps repository areas to responsibilities.

\begin{table*}[t]
\centering
\caption{Repository structure (selected).}
\label{tab:modules}
\begin{tabular}{@{}p{0.28\linewidth} p{0.65\linewidth}@{}}
\toprule
Path & Role \\
\midrule
\texttt{api/main.py}, \texttt{api/services/}, \texttt{api/agents/} & FastAPI lifespan, dataset bootstrap, multi-agent orchestration, chat/insert. \\
\texttt{model/} & \texttt{SocialGraphGNN}, inference manager, training utilities. \\
\texttt{db/}, \texttt{rag/} & Neo4j client; \texttt{hybrid\_retrieval.py} graph + vector retrieval; vector schema/population. \\
\texttt{training/} & Dataset-specific trainers (Facebook, Twitter, Reddit). \\
\texttt{tests/} & Unit/integration tests; \texttt{testing\_pipeline\_accuracy.py} for stratified pipeline accuracy. \\
\bottomrule
\end{tabular}
\end{table*}

\section{Training on Kaggle}

Training scripts run in GPU Kaggle notebooks: attach the appropriate dataset (e.g., Facebook MUSAE), write weights to \texttt{/kaggle/working}, then download \texttt{model\_weights\_*.pth} and embeddings into local \texttt{weights/}. Figures~\ref{fig:nb1}--\ref{fig:nb3} show the three evolution stages corresponding to \texttt{research\_paper\_implementation/} notebooks (RAG-only, GNN-only, GNN\,+\,RAG hybrid).

\begin{figure}[t]
\centering
\includegraphics[width=0.95\linewidth]{notebook1.png}
\caption{RAG-only exploration: embeddings and retrieval without graph structure.}
\label{fig:nb1}
\end{figure}

\begin{figure}[t]
\centering
\includegraphics[width=0.95\linewidth]{notebook2.png}
\caption{GNN-only baseline: GraphSAGE training and stratified evaluation.}
\label{fig:nb2}
\end{figure}

\begin{figure}[t]
\centering
\includegraphics[width=0.95\linewidth]{notebook3.png}
\caption{GNN\,+\,semantic feature hybrid training and checkpoint export.}
\label{fig:nb3}
\end{figure}

\section{Testing and Evaluation}

The suite under \texttt{tests/} includes feature and service tests (\texttt{test\_new\_features.py}, \texttt{test\_all.py}, etc.). The file \texttt{testing\_pipeline\_accuracy.py} uses the \texttt{testing\_*.py} naming convention; the repository includes \texttt{pytest.ini} so this file is discovered when running \texttt{pytest tests/}.

The pipeline accuracy test loads saved Facebook weights, builds a stratified train/validation/test split on labeled nodes, optionally fine-tunes the classifier with OneCycle learning-rate scheduling, constructs multi-hop propagation features, and selects blend/gate hyperparameters on the validation set only---then asserts a minimum test accuracy (configurable via environment variable). JSON metrics may be written under \texttt{/kaggle/working} when present.

\section{Results and Discussion}

Representative milestones from project documentation: GNN-only $\sim$93.5\% test accuracy; hybrid training $\sim$94.6\% test accuracy; the final inference-time graph-RAG pipeline targets at least $\sim$94\% on the held-out test slice with regression assertions~\cite{project_doc_internal}. Exact numbers depend on hardware, stochasticity, and checkpoint version.

\section{Conclusion}

We described a GraphRAG-style social intelligence stack integrating GNNs, Neo4j-backed hybrid retrieval, and LLM synthesis, grounded in survey literature on graph ML with LLMs~\cite{acm_survey_gml_llm}. Future work includes scaling to larger graphs, latency optimization for vector search, and extended evaluation on Twitter and Reddit datasets already scaffolded in the repository.

\section*{Acknowledgment}

The authors thank Shiv Nadar University and the Department of Computer Science for resources and support for this course project.

\bibliographystyle{IEEEtran}
\begin{thebibliography}{00}

\bibitem{acm_survey_gml_llm}
``Graph Machine Learning in the Era of Large Language Models,'' \emph{ACM Comput. Surv.}, 2025.

\bibitem{project_doc_internal}
``Facebook Graph Classification --- Project Documentation,'' training log, GraphRAG Social Intelligence repository, \texttt{training/project\_documentation.md}.

\bibitem{arch_md_internal}
``Neo4j-Only Architecture: Design Decisions \& Performance Analysis,'' GraphRAG Social Intelligence repository, \texttt{ARCHITECTURE.md}.

\bibitem{hamilton2017graphsage}
W.~L. Hamilton, R.~Ying, and J.~Leskovec, ``Inductive representation learning on large graphs,'' in \emph{Proc. NeurIPS}, 2017.

\bibitem{pyg}
M.~Fey and J.~E. Lenssen, ``Fast graph representation learning with PyTorch Geometric,'' in \emph{ICLR Workshop on Representation Learning on Graphs and Manifolds}, 2019.

\bibitem{neo4j_vector}
Neo4j, ``Vector indexes,'' documentation, \url{https://neo4j.com/docs/cypher-manual/current/indexes-for-vector-search/}.

\end{thebibliography}

\end{document}
